\begin{description}
\item[1)] The luminosity distance is defined to satisfy a canonical inverse
  square law of the bolometric luminosity $L$ (total emitted radiation power)
  and observed flux $F$ (flux--luminosity relationship)
  \[ F = \frac{L}{4\pi d_L^2} \tag*{(1)} \]
  Units of the flux and luminosity are expressed in SI
  ($\mathrm{W\,m^{-2}}$ and W respectively), or CGS
  ($\mathrm{erg\,s^{-1}\,cm^{-2}}$, $\mathrm{erg\,s^{-1}}$ respectively).
  \hfill(3 points)

\item[2)] Only the part of the spectrum (namely B-band) is observed, but the
  inverse square law is applicable also to the monochromatic luminosity
  $L_\lambda(\lambda)$ and monochromatic flux
  $S(\lambda)$,\footnote{Monochromatic flux is sometimes called as spectral
  density of flux (energy per unit time per unit area per unit wavelength)
  while monochromatic luminosity as spectral density of the luminosity (energy
  per unit time per unit wavelength). Their form is described by SED and thus
  spectral index.} taking into account that observed wavelengths are not equal
  to emitted ones. (Observers frame is not the same as a frame of galaxy.) For
  a small wavelength interval
  \[ S(\lambda_o)\cdot\Delta\lambda_o
     = \frac{L_\lambda(\lambda_e)}{4\pi d_L^2}\cdot\Delta\lambda_e , \]
  Observed fraction of the total spectrum is redshifted into the frame of the
  observer, and a relationship between the observed and emitted wavelengths,
  $\lambda_o$ and $\lambda_e$, can be expressed as
  \[ \lambda_e = \lambda_o/(1+z) . \]
  So,
  \[ S(\lambda)\Delta\lambda = \frac{1}{4\pi d_L^2}\cdot
     L_\lambda\!\left(\frac{\lambda}{1+z}\right)\cdot\frac{1}{1+z}\,
     \Delta\lambda . \tag*{(2)} \]
  \hfill(9 points)

\item[3)] The absolute luminosity is defined as the flux would be observed at
  a distance of 10\,pc from the compact source at rest (not redshifted). It
  means that the frame would be the same as for the object.
  \[ S_{10}(\lambda) = \frac{L_\lambda(\lambda)}{4\pi d_{10}^2}, \tag*{(3)} \]
  where $d_{10}$ is the distance equal to 10\,pc, $S_{10}(\lambda)$ denotes a
  monochromatic flux in the wavelength of $\lambda$ measured at 10\,pc from
  the object.
  \hfill(3 points)

\item[4)] Dividing eq.~(2) by (3) we get
  \[ \frac{S(\lambda)}{S_{10}(\lambda)}
     = \left(\frac{d_{10}}{d_L}\right)^2 \cdot
       \frac{L_\lambda\!\left(\dfrac{\lambda}{1+z}\right)}{L_\lambda(\lambda)}
       \cdot \frac{1}{1+z} \tag*{(4)} \]
  for all the wavelengths.

  The lower edge of validity of SED approximation is equal to 250\,nm. This
  corresponds to $1.5\cdot 250\ \mathrm{nm} = 375$\,nm what is below the blue
  cut-off of the B filter,\footnote{Blue band: Effective Wavelength Midpoint
  $\lambda\mathrm{eff}=445$\,nm, FWHM${}=94$\,nm
  [398\,nm $\div$ 492\,nm]} so SED approximation is applicable. Using the SED
  of the galaxy, we get
  \begin{align*}
    \frac{L_\lambda\!\left(\dfrac{\lambda}{1+z}\right)}{L_\lambda(\lambda)}
      \cdot \frac{1}{1+z}
      &= \left(\frac{1}{1+z}\right)^{4}\frac{1}{1+z}
       = (1+z)^{-5} \tag*{(5)} \\
    \frac{S(\lambda)}{S_{10}(\lambda)}
      &= \left(\frac{d_{10}}{d_L}\right)^{2}\cdot(1+z)^{-5} \tag*{(6)}
  \end{align*}
  The fraction of the monochromatic fluxes $S(\lambda)/S_{10}(\lambda)$ does
  not depend on the wavelength! So the relation (6) is valid also for blue
  band fluxes.
  \hfill(9 points)

\item[5)] The blue fluxes $F$ and $F_{10}$ have to be replaced by the
  corresponding magnitudes. Using Pogson's Ratio it can be obtained:
  \begin{align*}
    m_B - M_B &= -2.5\log\left(\frac{F}{F_{10}}\right)
      = 5\log\left(\frac{d_L}{d_{10}}\right) + 12.5\log(1+z) \\
    M_B &= m_B - 5\log\left(\frac{d_L}{d_{10}}\right)
      + 12.5\log(1+z) \tag*{(7)} % sic: source has "+"; the numerical
      % substitution below (and the first line) require "-"
  \end{align*}
  Substituting numerical data one can obtain:
  \[ M_B = 20.40^{\mathrm{mag}} - 5\log\!\left(2754\cdot 10^{5}\right)
     - 12.5\log(1.5) = -24.00^{\mathrm{mag}} \]
  \[ M_B = -24.00^{\mathrm{mag}} \tag*{(8)} \]
  Please note, that precision is determined by $m_B$!
  \hfill(3 points)

\item[6)] Only the cD galaxies are so luminous, while the normal elliptical
  galaxies are substantially weaker. Accordingly, the galaxy is not a member
  of the cluster, but is a foreground object.
  \hfill(3 points)
\end{description}
